% =============================================================================
% report-card template — per-student grade reports
% =============================================================================
%
% Compile with:
%   lualatex report-card-template.tex

% Course-level fields are scalars, so they ride in the class option list.
% (Real card sets usually inherit these from a coursemeta.tex up the tree.)
\documentclass[
	report-date    = {May 7, 2026},
	institution    = {University of Nevada, Reno},
	instructor     = {Your Name},
	season         = Spring,
	year           = 2026,
	course-subject = Math,
	course-number  = 181,
	course-title   = {Calculus I},
	gradebook-path = {gradebook.xlsx},
]{report-card}

\begin{document}

% One report card per student, each with its own cover. Grades — and all the
% arithmetic — live in the gradebook sheet; gradebook.csv is its report-view
% export (the only tab this document reads). Columns follow the naming
% convention "<Category> Weight/Score/Points", "Current Total/Points", and
% "Need <letter>"; a "---" column inserts a midrule in the breakdown table.
%
% Bare \gradebook resolves the path from coursemeta's gradebook-path above
% (here gradebook.xlsx, mapped to the gradebook.csv the builder exports),
% falling back to a sibling gradebook.csv. Pass \gradebook[other.csv] to
% override.
\gradebook

\end{document}

% ----------------------------------------------------------------------------
% Manual escape hatch — type a single card by hand without the gradebook:
%
% \begin{reportcard}{Another Student}
%     \section*{Current Grade Calculation}
%     \gradebreakdown{
%         Homework Avg. & 15\% & 85.0\% & +12.8 \\
%         \midrule
%         Exam 1        & 10\% & 78\%   & +7.8  \\
%     }{75\% (+15\% E.C.) & \fbox{\textbf{86.0\%}} & 64.5}
%     \standingbar{10.75}{7.5}{8.75}{10}{11.25}
%     \scenarios{ To earn a C & C 70\% & \textbf{22.0\%} \\ }
%     \signoff{Thank you for attending my course!}
% \end{reportcard}
